\usepackage{geometry}
\geometry{
  a4paper,
  total={160mm,257mm},
  left=25mm,
  top=20mm,
}

\usepackage{stmaryrd}
\usepackage{bussproofs}
\usepackage{listings}
\usepackage[utf8]{inputenc}
\usepackage{fancyvrb}
\usepackage{amsmath}
\usepackage{amsfonts}
\usepackage{amssymb}
\usepackage{graphicx}
\usepackage{subcaption}
\usepackage{hyperref}
\usepackage{sectsty}
\usepackage{titlesec}
\usepackage{amssymb}
\usepackage[dvipsnames,x11names]{xcolor}
\usepackage{tikz}
\usepackage{tikz-cd}
\usepackage{wrapfig}
\usepackage{array}
\usepackage{dirtytalk}
\usepackage{ifthen}
\usepackage{framed}
\usepackage{tcolorbox}
\usepackage{unicode-math}
\usepackage{pgfpages}
\usepackage{pgfkeys}
\usepackage{setspace}
\usepackage{iftex}
\usepackage{keyval}
\usepackage{mdframed}
\usepackage{biblatex} %Imports biblatex package

\usepackage{iftex}
\ifPDFTeX
  \usepackage[T1]{fontenc}
  \usepackage[utf8]{inputenc}
  \usepackage{textcomp} % provide euro and other symbols
\else % if luatex or xetex
  \usepackage{unicode-math} % this also loads fontspec
  % \defaultfontfeatures{Scale=MatchLowercase}
  % \defaultfontfeatures[\rmfamily]{Ligatures=TeX,Scale=1}
\fi
\usepackage{lmodern}
\ifPDFTeX\else
  % xetex/luatex font selection
\fi
% Use upquote if available, for straight quotes in verbatim environments
\IfFileExists{upquote.sty}{\usepackage{upquote}}{}
\IfFileExists{microtype.sty}{% use microtype if available
  \usepackage[]{microtype}
  \UseMicrotypeSet[protrusion]{basicmath} % disable protrusion for tt fonts
}{}
\makeatletter
\@ifundefined{KOMAClassName}{% if non-KOMA class
  \IfFileExists{parskip.sty}{%
    \usepackage{parskip}
  }{% else
    \setlength{\parindent}{0pt}
    \setlength{\parskip}{6pt plus 2pt minus 1pt}}
}{% if KOMA class
  \KOMAoptions{parskip=half}}
\makeatother
\usepackage{color}
\usepackage{fancyvrb}
\newcommand{\VerbBar}{|}
\newcommand{\VERB}{\Verb[commandchars=\\\{\}]}
\DefineVerbatimEnvironment{Highlighting}{Verbatim}{
    commandchars=\\\{\},
    frame=lines,
    rulecolor=\color{lightgray}}
% Add ',fontsize=\small' for more characters per line
\newenvironment{Shaded}{}{}
\newcommand{\AlertTok}[1]{\textcolor[rgb]{1.00,0.00,0.00}{\textbf{#1}}}
\newcommand{\AnnotationTok}[1]{\textcolor[rgb]{0.38,0.63,0.69}{\textbf{\textit{#1}}}}
\newcommand{\AttributeTok}[1]{\textcolor[rgb]{0.49,0.56,0.16}{#1}}
\newcommand{\BaseNTok}[1]{\textcolor[rgb]{0.25,0.63,0.44}{#1}}
\newcommand{\BuiltInTok}[1]{\textcolor[rgb]{0.00,0.50,0.00}{#1}}
\newcommand{\CharTok}[1]{\textcolor[rgb]{0.25,0.44,0.63}{#1}}
\newcommand{\CommentTok}[1]{\textcolor[rgb]{0.38,0.63,0.69}{\textit{#1}}}
\newcommand{\CommentVarTok}[1]{\textcolor[rgb]{0.38,0.63,0.69}{\textbf{\textit{#1}}}}
\newcommand{\ConstantTok}[1]{\textcolor[rgb]{0.53,0.00,0.00}{#1}}
\newcommand{\ControlFlowTok}[1]{\textcolor[rgb]{0.00,0.44,0.13}{\textbf{#1}}}
\newcommand{\DataTypeTok}[1]{\textcolor[rgb]{0.56,0.13,0.00}{#1}}
\newcommand{\DecValTok}[1]{\textcolor[rgb]{0.25,0.63,0.44}{#1}}
\newcommand{\DocumentationTok}[1]{\textcolor[rgb]{0.73,0.13,0.13}{\textit{#1}}}
\newcommand{\ErrorTok}[1]{\textcolor[rgb]{1.00,0.00,0.00}{\textbf{#1}}}
\newcommand{\ExtensionTok}[1]{#1}
\newcommand{\FloatTok}[1]{\textcolor[rgb]{0.25,0.63,0.44}{#1}}
\newcommand{\FunctionTok}[1]{\textcolor[rgb]{0.02,0.16,0.49}{#1}}
\newcommand{\ImportTok}[1]{\textcolor[rgb]{0.00,0.50,0.00}{\textbf{#1}}}
\newcommand{\InformationTok}[1]{\textcolor[rgb]{0.38,0.63,0.69}{\textbf{\textit{#1}}}}
\newcommand{\KeywordTok}[1]{\textcolor[rgb]{0.00,0.44,0.13}{\textbf{#1}}}
\newcommand{\NormalTok}[1]{#1}
\newcommand{\OperatorTok}[1]{\textcolor[rgb]{0.40,0.40,0.40}{#1}}
\newcommand{\OtherTok}[1]{\textcolor[rgb]{0.00,0.44,0.13}{#1}}
\newcommand{\PreprocessorTok}[1]{\textcolor[rgb]{0.74,0.48,0.00}{#1}}
\newcommand{\RegionMarkerTok}[1]{#1}
\newcommand{\SpecialCharTok}[1]{\textcolor[rgb]{0.25,0.44,0.63}{#1}}
\newcommand{\SpecialStringTok}[1]{\textcolor[rgb]{0.73,0.40,0.53}{#1}}
\newcommand{\StringTok}[1]{\textcolor[rgb]{0.25,0.44,0.63}{#1}}
\newcommand{\VariableTok}[1]{\textcolor[rgb]{0.10,0.09,0.49}{#1}}
\newcommand{\VerbatimStringTok}[1]{\textcolor[rgb]{0.25,0.44,0.63}{#1}}
\newcommand{\WarningTok}[1]{\textcolor[rgb]{0.38,0.63,0.69}{\textbf{\textit{#1}}}}
\setlength{\emergencystretch}{3em} % prevent overfull lines
\providecommand{\tightlist}{%
  \setlength{\itemsep}{0pt}\setlength{\parskip}{0pt}}
\setcounter{secnumdepth}{-\maxdimen} % remove section numbering
\usepackage{fancyvrb}
\usepackage[x11names]{xcolor}
\newcommand{\Katla}                [2][]{\VerbatimInput[commandchars=\\\{\}#1]{#2}}
\newcommand{\KatlaNewline}            {}
\newcommand{\KatlaSpace}              { }
\newcommand{\KatlaDash}               {\string-}
\newcommand{\KatlaUnderscore}         {\string_}
\newcommand{\KatlaTilde}              {\string~}
\newcommand{\IdrisHlightFont}         {\ttfamily}
\newcommand{\IdrisHlightStyleData}    {}
\newcommand{\IdrisHlightStyleType}    {}
\newcommand{\IdrisHlightStyleBound}   {}
\newcommand{\IdrisHlightStyleFunction}{}
\newcommand{\IdrisHlightStyleKeyword} {\bfseries}
\newcommand{\IdrisHlightStyleImplicit}{}
\newcommand{\IdrisHlightStyleComment} {\itshape}
\newcommand{\IdrisHlightStyleHole}    {\bfseries}
\newcommand{\IdrisHlightStyleNamespace}{\itshape}
\newcommand{\IdrisHlightStylePostulate}{\bfseries}
\newcommand{\IdrisHlightStyleModule}   {\itshape}

\newcommand{\IdrisHlightColourData}    {IndianRed1}
\newcommand{\IdrisHlightColourType}    {DeepSkyBlue3}
\newcommand{\IdrisHlightColourBound}   {DarkOrchid3}
\newcommand{\IdrisHlightColourFunction}{Chartreuse4}
\newcommand{\IdrisHlightColourKeyword} {black}
\newcommand{\IdrisHlightColourImplicit}{DarkOrchid3}
\newcommand{\IdrisHlightColourComment} {Ivory3}
\newcommand{\IdrisHlightColourHole}    {yellow}
\newcommand{\IdrisHlightColourNamespace}{black}
\newcommand{\IdrisHlightColourPostulate}{DarkOrchid3}
\newcommand{\IdrisHlightColourModule}   {black}

\newcommand{\IdrisHole}[1]{{%
    \colorbox{\IdrisHlightColourHole}{%
      \IdrisHlightStyleHole\IdrisHlightFont%
      #1}}}

\newcommand{\RawIdrisHighlight}[3]{{\textcolor{#1}{\IdrisHlightFont#2{#3}}}}

\newcommand{\IdrisData}[1]{\RawIdrisHighlight{\IdrisHlightColourData}{\IdrisHlightStyleData}{#1}}
\newcommand{\IdrisType}[1]{\RawIdrisHighlight{\IdrisHlightColourType}{\IdrisHlightStyleType}{#1}}
\newcommand{\IdrisBound}[1]{\RawIdrisHighlight{\IdrisHlightColourBound}{\IdrisHlightStyleBound}{#1}}
\newcommand{\IdrisFunction}[1]{\RawIdrisHighlight{\IdrisHlightColourFunction}{\IdrisHlightStyleFunction}{#1}}
\newcommand{\IdrisKeyword}[1]{\RawIdrisHighlight{\IdrisHlightColourKeyword}{\IdrisHlightStyleKeyword}{#1}}
\newcommand{\IdrisImplicit}[1]{\RawIdrisHighlight{\IdrisHlightColourImplicit}{\IdrisHlightStyleImplicit}{#1}}
\newcommand{\IdrisComment}[1]{\RawIdrisHighlight{\IdrisHlightColourComment}{\IdrisHlightStyleComment}{#1}}
\newcommand{\IdrisNamespace}[1]{\RawIdrisHighlight{\IdrisHlightColourNamespace}{\IdrisHlightStyleNamespace}{#1}}
\newcommand{\IdrisPostulate}[1]{\RawIdrisHighlight{\IdrisHlightColourPostulate}{\IdrisHlightStylePostulate}{#1}}
\newcommand{\IdrisModule}[1]{\RawIdrisHighlight{\IdrisHlightColourModule}{\IdrisHlightStyleModule}{#1}}

\DefineVerbatimEnvironment%
{code}{Verbatim}{commandchars=\\\{\}, frame=single}

% Bugfix in fancyvrb to allow inline saved listings
\makeatletter
\let\FV@ProcessLine\relax
\makeatother
\ifLuaTeX
  \usepackage{selnolig}  % disable illegal ligatures
\fi
\usepackage{bookmark}
\IfFileExists{xurl.sty}{\usepackage{xurl}}{} % add URL line breaks if available
\urlstyle{same}
\hypersetup{
  hidelinks,
  pdfcreator={LaTeX via pandoc}}

\usepackage{fontspec}

% Simple active character approach
% \catcode`⊗=\active
% \protected\def⊗{⊗} % This preserves the character as-is
\setmonofont[
    Path = ../font/JetBrains_Mono/static/,
    Renderer=Harfbuzz,
    UprightFont = *-Regular,
    ItalicFont = *-Italic,
    Scale=0.9]
    {JetBrainsMono}
\providecommand{\tightlist}{%
  \setlength{\itemsep}{0pt}\setlength{\parskip}{0pt}
}

\usepackage{pmboxdraw}
\usepackage{newunicodechar}
%\newunicodechar{└}{\textSFii}
%\newunicodechar{├}{\textSFviii}
%\newunicodechar{─}{\textSFx}
%\newunicodechar{∀}{$\forall$}
%\newunicodechar{η}{$\eta$}
%\newunicodechar{μ}{$\mu$}
%\newunicodechar{→}{$\to$}
%\newunicodechar{⊕}{$\oplus$}
%\newunicodechar{∈}{$\in$}
%\newunicodechar{∋}{$\ni$}
%\newunicodechar{≡}{$\equiv$}
%\newunicodechar{⊃}{$\supset$}
%\newunicodechar{⊂}{$\subset$}
%\newunicodechar{√}{$\sqrt{}$}
%\newunicodechar{∫}{$\int$}
%\newunicodechar{╲}{\textbackslash} % Not sure about this one
%\newunicodechar{○}{$\circ$}
%\newunicodechar{Σ}{$\Sigma$}
%\newunicodechar{π}{$\pi$}
%\newunicodechar{λ}{$\lambda$}
\newunicodechar{⨾}{;}
%\newunicodechar{⊗}{$\otimes$}
%% For XeLaTeX only
%\XeTeXinterchartokenstate=1
%\newXeTeXintercharclass\unicodemath
%\XeTeXcharclass`⊗=\unicodemath
%\XeTeXinterchartoks 0 \unicodemath = {\ensuremath{\otimes}}
%
%\newunicodechar{Π}{$\Pi$}
%\newunicodechar{φ}{$\varphi$}

\newcommand{\of}{\text{ }}
\newcommand{\cat}[1]{\mathcal{#1}}
\newcommand{\Poly}{\cat{Poly}}
\newcommand{\Cont}{\cat{Cont}}
\newcommand{\Set}{\cat{Set}}
\newcommand{\ContCart}{\cat{Cont}^{\#}}
\newcommand{\Type}{\mathit{Type}}
\newcommand{\mor}[3]{\mathcal{#1}(#2, #3)}  % morphisms
\newcommand{\cbar}[1]{(#1, \bar{#1})} % (A, \bar A) containers
% The traditional composition
\newcommand{\compose}{\mathbin{\rhd}}
% The "universal composition" monoidal action
\newcommand{\forallSeq}{\mathbin{\blacktriangleright}}
\newcommand{\UnivComp}{\forallSeq}

\newcommand{\figref}[1]{(\textbf{Figure}~\ref{#1})}
\newcommand{\secref}[1]{(\textbf{Section}~\ref{#1})}
\newcommand{\defref}[1]{(\textbf{Definition}~\ref{#1})}
\newcommand{\lemref}[1]{(\textbf{Lemma}~\ref{#1})}
\newcommand{\propref}[1]{(\textbf{Proposition}~\ref{#1})}
\newcommand{\thmref}[1]{(\textbf{Theorem}~\ref{#1})}
\newcommand{\postref}[1]{(\textbf{Postulate}~\ref{#1})}
\newcommand{\comp}{\mathbin{;}}
\newcommand{\ncomp}[1]{\mathbin{;_#1}}
\newcommand{\vcomp}{\mathbin{;_v}}
\newcommand{\hcomp}{\mathbin{;_h}}


\usetikzlibrary{graphs}
\usetikzlibrary{arrows.meta}
\makeatletter

\newcommand*{\RightArrow}{\mathrel{\mathpalette{\@RightArrow}{}}}
\newcommand*{\@RightArrow}[1]{%
    % Get the line width for this math style
    \edef\@LineWidth{%
        \the\fontdimen8
        \ifx#1\displaystyle\textfont
        \else\ifx#1\textstyle\textfont
        \else\ifx#1\scriptstyle\scriptfont
        \else\scriptscriptfont
        \fi\fi\fi
        3}
    \edef\@ScaleWidth{%
        \ifx#1\displaystyle0.39
        \else\ifx#1\textstyle0.39
        \else\ifx#1\scriptstyle0.35
        \else0.31
        \fi\fi\fi}
    \text{$\tikz
    \draw[double equal sign distance, line width=\@LineWidth,
    -{Implies[sep=-0.5ex] . Computer Modern Rightarrow[scale width=\@ScaleWidth, scale length=0.8]}]
    (0,0) -- (1.5em,0);$}}
\newcommand*{\Implies}{\DOTSB\;\RightArrow\;}

% Theorems, declarations, etc
% Proposition colors
\mdfdefinestyle{postulateframe}{
    leftmargin=0pt,
    rightmargin=0pt,
    innerleftmargin=15pt,
    innerrightmargin=10pt,
    innertopmargin=8pt,
    innerbottommargin=8pt,
    linewidth=4pt,
    topline=false,
    rightline=false,
    bottomline=false,
    leftline=true,
    linecolor=red!30,
    backgroundcolor=red!3
}

% Theorem colors
\mdfdefinestyle{theoremframe}{
    leftmargin=0pt,
    rightmargin=0pt,
    innerleftmargin=15pt,
    innerrightmargin=10pt,
    innertopmargin=8pt,
    innerbottommargin=8pt,
    linewidth=4pt,
    topline=false,
    rightline=false,
    bottomline=false,
    leftline=true,
    linecolor=blue,
    linecolor=purple!30,
    backgroundcolor=purple!3
}
% Theorem colors
\mdfdefinestyle{lemmaframe}{
    leftmargin=0pt,
    rightmargin=0pt,
    innerleftmargin=15pt,
    innerrightmargin=10pt,
    innertopmargin=8pt,
    innerbottommargin=8pt,
    linewidth=4pt,
    topline=false,
    rightline=false,
    bottomline=false,
    leftline=true,
    linecolor=yellow!20,
    backgroundcolor=yellow!1
}

% Theorem colors
\mdfdefinestyle{propositionframe}{
    leftmargin=0pt,
    rightmargin=0pt,
    innerleftmargin=15pt,
    innerrightmargin=10pt,
    innertopmargin=8pt,
    innerbottommargin=8pt,
    linewidth=4pt,
    topline=false,
    rightline=false,
    bottomline=false,
    leftline=true,
    %jlinecolor=blue,
    linecolor=orange!30,
    backgroundcolor=orange!2
}
% Theorem colors
\mdfdefinestyle{definitionframe}{
    leftmargin=0pt,
    rightmargin=0pt,
    innerleftmargin=15pt,
    innerrightmargin=10pt,
    innertopmargin=8pt,
    innerbottommargin=8pt,
    linewidth=4pt,
    topline=false,
    rightline=false,
    bottomline=false,
    leftline=true,
    linecolor=green!30,
    backgroundcolor=green!2
}

\newtheorem{prop}{Proposition}
\newtheorem{thm}{Theorem}
\newtheorem{lemma}{Lemma}
\newtheorem{defn}{Definition}
\newtheorem{post}{Postulate}


% Apply the frame to your existing theorem
\surroundwithmdframed[style=definitionframe]{defn}
\surroundwithmdframed[style=propositionframe]{prop}
\surroundwithmdframed[style=lemmaframe]{lemma}
\surroundwithmdframed[style=theoremframe]{thm}
\surroundwithmdframed[style=postulateframe]{post}

%% Pentagon diagram macro
\pgfkeys{
  /pentagon/.cd,
  node top left/.code={\def\penttopleft{#1}},
  node top right/.code={\def\penttopright{#1}},
  node mid left/.code={\def\pentmidleft{#1}},
  node bottom left/.code={\def\pentbottomleft{#1}},
  node bottom right/.code={\def\pentbottomright{#1}},
  arrow top/.code={\def\penttop{#1}},
  arrow left/.code={\def\pentleft{#1}},
  arrow right/.code={\def\pentright{#1}},
  arrow bottom left/.code={\def\pentbottomleftarrow{#1}},
  arrow bottom right/.code={\def\pentbottomrightarrow{#1}},
}

% Initialize all macros to avoid undefined errors
\def\penttopleft{}
\def\penttopright{}
\def\pentmidleft{}
\def\pentbottomleft{}
\def\pentbottomright{}
\def\penttop{}
\def\pentleft{}
\def\pentright{}
\def\pentbottomleftarrow{}
\def\pentbottomrightarrow{}

\newcommand{\pentagondiagram}[1]{%
  \pgfkeys{/pentagon/.cd, #1}%
  \begin{tikzcd}[ampersand replacement=\&]
    {\penttopleft} \& \& {\penttopright} \\
    {\pentmidleft} \& \& {} \\
    {\pentbottomleft} \& \& {\pentbottomright}
    \arrow[from=1-1, to=1-3, "\penttop"]
    \arrow[from=1-1, to=2-1, "\pentleft"]
    \arrow[from=1-3, to=3-3, "\pentright"]
    \arrow[from=2-1, to=3-1, "\pentbottomleftarrow"]
    \arrow[from=3-1, to=3-3, "\pentbottomrightarrow"]
  \end{tikzcd}
}

\addbibresource{Thesis-bib.bib} %Import the bibliography file

